\\documentclass[11pt]{article}
\\usepackage[T1]{fontenc}
\\usepackage[utf8]{inputenc}
\\usepackage{geometry}
\\usepackage{xcolor}
\\usepackage{fvextra}
\\usepackage{hyperref}
\\geometry{margin=1in}
\\DefineVerbatimEnvironment{CodeBlock}{Verbatim}{breaklines,breakanywhere,fontsize=\\small}

\begin{document}
\section*{core/rks\\_diagnostic\\_experiment.py}
\textbf{Source Path}: \texttt{core/rks\\_diagnostic\\_experiment.py}\par
\bigskip
\begin{CodeBlock}
#!/usr/bin/env python3
"""
rks_diagnostic_experiment.py

THE SINGLE MOST IMPORTANT EXPERIMENT for validating your thesis.

This script answers:
1. Is RFF approximation good enough? (correlation > 0.95)
2. Do "observer residuals" shrink as D → ∞? (if yes, they're Monte Carlo noise)
3. Is your geometry sensitive to sigma? (if highly sensitive, it's an artifact)

If residuals persist as D grows AND approximation is good, you have a real phenomenon.
If residuals vanish as D grows, seeds were never observers — just Monte Carlo noise.

Usage:
    python rks_diagnostic_experiment.py --data-path experiments/rbf/logits/real/
    python rks_diagnostic_experiment.py --synthetic --input-dim 24 --n-samples 500
"""

import argparse
import json
import torch
import numpy as np
from pathlib import Path
from typing import Dict, List, Optional
from dataclasses import dataclass, asdict
import matplotlib.pyplot as plt

# Handle imports
try:
    from kernel_library import create_kernel, SUPPORTED_KERNELS, run_approximation_diagnostics
except ImportError:
    import sys
    sys.path.insert(0, str(Path(__file__).parent))
    from kernel_library import create_kernel, SUPPORTED_KERNELS, run_approximation_diagnostics


@dataclass
class DiagnosticResult:
    kernel_type: str
    output_dim: int
    sigma: float
    correlation: float
    mae: float
    rmse: float
    
    def is_acceptable(self, threshold: float = 0.95) -> bool:
        return self.correlation >= threshold


def load_nli_features(data_path: Path) -> torch.Tensor:
    """Load NLI features from observer files."""
    files = sorted(data_path.glob("observer_*.pt"))
    if not files:
        files = sorted(data_path.glob("seed_*.pt"))
    
    if not files:
        raise FileNotFoundError(f"No observer files in {data_path}")
    
    # Load first file to get features
    data = torch.load(files[0], map_location='cpu', weights_only=False)
    
    if isinstance(data, dict):
        for key in ['features', 'embeddings', 'nli_features']:
            if key in data:
                return torch.FloatTensor(data[key])
    
    if torch.is_tensor(data):
        return data.float()
    
    raise ValueError(f"Could not extract features from {files[0]}")


def compute_observer_residuals(
    X: torch.Tensor,
    kernel_type: str,
    output_dim: int,
    sigma: float,
    seeds: List[int],
) -> Dict:
    """
    Compute residuals between observers (different seeds).
    
    This measures: how much does geometry change with seed?
    """
    from scipy.linalg import orthogonal_procrustes
    
    embeddings = []
    for seed in seeds:
        kernel = create_kernel(kernel_type, X.shape[1], output_dim, sigma, seed)
        phi = kernel.transform(X)
        embeddings.append(phi.cpu().numpy())
    
    # Use first observer as reference
    ref = embeddings[0]
    
    # Align others to reference via Procrustes
    residuals = []
    for i, emb in enumerate(embeddings[1:], 1):
        R, _ = orthogonal_procrustes(emb, ref)
        aligned = emb @ R
        
        # Frobenius distance after alignment
        residual = np.linalg.norm(aligned - ref, ord='fro')
        residuals.append(residual)
    
    return {
        'mean_residual': float(np.mean(residuals)),
        'std_residual': float(np.std(residuals)),
        'residuals': residuals,
        'n_observers': len(seeds),
    }


def run_d_sweep_experiment(
    X: torch.Tensor,
    kernel_type: str = 'rbf',
    output_dims: List[int] = None,
    sigma: float = None,
    seeds: List[int] = None,
) -> Dict:
    """
    THE KEY EXPERIMENT: How do observer residuals change with D?
    
    If residuals → 0 as D → ∞: seeds are just Monte Carlo noise
    If residuals persist: there might be real observer structure
    """
    if output_dims is None:
        output_dims = [256, 512, 1024, 2048, 4096]
    
    if seeds is None:
        seeds = list(range(42, 52))  # 10 observers
    
    if sigma is None:
        dists = torch.cdist(X[:100], X[:100])
        sigma = float(dists.median())
    
    results = {
        'kernel_type': kernel_type,
        'sigma': sigma,
        'n_samples': len(X),
        'input_dim': X.shape[1],
        'seeds': seeds,
        'by_dim': {},
    }
    
    print(f"\n{'='*60}")
    print(f"D-SWEEP EXPERIMENT: {kernel_type}, σ={sigma:.4f}")
    print(f"{'='*60}")
    print(f"{'D':<8} {'Corr':<10} {'Residual':<12} {'Residual/D':<12}")
    print("-" * 45)
    
    for D in output_dims:
        # Approximation quality
        kernel = create_kernel(kernel_type, X.shape[1], D, sigma, seed=42)
        diag = kernel.approximation_diagnostic(X)
        
        # Observer residuals
        resid = compute_observer_residuals(X, kernel_type, D, sigma, seeds)
        
        # Normalized residual (should → 0 if just Monte Carlo noise)
        normalized = resid['mean_residual'] / np.sqrt(D)
        
        results['by_dim'][D] = {
            'correlation': diag['correlation'],
            'mae': diag['mae'],
            'rmse': diag['rmse'],
            'mean_residual': resid['mean_residual'],
            'std_residual': resid['std_residual'],
            'normalized_residual': normalized,
        }
        
        print(f"{D:<8} {diag['correlation']:<10.4f} {resid['mean_residual']:<12.4f} {normalized:<12.6f}")
    
    # Analysis: do residuals shrink with D?
    dims = sorted(results['by_dim'].keys())
    residuals = [results['by_dim'][d]['mean_residual'] for d in dims]
    normalized = [results['by_dim'][d]['normalized_residual'] for d in dims]
    
    # Fit power law: residual ~ D^(-alpha)
    log_d = np.log(dims)
    log_r = np.log(residuals)
    slope, intercept = np.polyfit(log_d, log_r, 1)
    
    results['analysis'] = {
        'power_law_exponent': float(slope),
        'residuals_shrinking': slope < -0.3,  # Should be ~ -0.5 for pure Monte Carlo
        'interpretation': (
            "MONTE CARLO NOISE: Residuals shrink as D^{:.2f}".format(slope)
            if slope < -0.3 else
            "POSSIBLE SIGNAL: Residuals persist (exponent = {:.2f})".format(slope)
        ),
    }
    
    print(f"\n{'='*60}")
    print(f"ANALYSIS: Residual ~ D^{slope:.3f}")
    print(f"{'='*60}")
    print(results['analysis']['interpretation'])
    
    return results


def run_sigma_sweep_experiment(
    X: torch.Tensor,
    kernel_type: str = 'rbf',
    output_dim: int = 1024,
    sigma_multipliers: List[float] = None,
    seeds: List[int] = None,
) -> Dict:
    """
    Test sensitivity to sigma choice.
    
    If geometry changes dramatically with sigma, it's an artifact.
    """
    if sigma_multipliers is None:
        sigma_multipliers = [0.1, 0.25, 0.5, 1.0, 2.0, 4.0, 10.0]
    
    if seeds is None:
        seeds = list(range(42, 52))
    
    # Base sigma from data
    dists = torch.cdist(X[:100], X[:100])
    base_sigma = float(dists.median())
    
    results = {
        'kernel_type': kernel_type,
        'output_dim': output_dim,
        'base_sigma': base_sigma,
        'by_sigma': {},
    }
    
    print(f"\n{'='*60}")
    print(f"SIGMA SWEEP: {kernel_type}, D={output_dim}")
    print(f"{'='*60}")
    print(f"{'σ mult':<10} {'σ':<12} {'Corr':<10} {'Residual':<12}")
    print("-" * 45)
    
    for mult in sigma_multipliers:
        sigma = base_sigma * mult
        
        kernel = create_kernel(kernel_type, X.shape[1], output_dim, sigma, seed=42)
        diag = kernel.approximation_diagnostic(X)
        resid = compute_observer_residuals(X, kernel_type, output_dim, sigma, seeds)
        
        results['by_sigma'][mult] = {
            'sigma': sigma,
            'correlation': diag['correlation'],
            'mean_residual': resid['mean_residual'],
        }
        
        print(f"{mult:<10.2f} {sigma:<12.4f} {diag['correlation']:<10.4f} {resid['mean_residual']:<12.4f}")
    
    # Check sensitivity
    residuals = [results['by_sigma'][m]['mean_residual'] for m in sigma_multipliers]
    cv = np.std(residuals) / np.mean(residuals)  # Coefficient of variation
    
    results['analysis'] = {
        'residual_cv': float(cv),
        'highly_sensitive': cv > 0.5,
        'interpretation': (
            "WARNING: Geometry highly sensitive to σ (CV={:.2f})".format(cv)
            if cv > 0.5 else
            "OK: Geometry stable across σ range (CV={:.2f})".format(cv)
        ),
    }
    
    print(f"\n{results['analysis']['interpretation']}")
    
    return results


def run_kernel_discrimination_experiment(
    X: torch.Tensor,
    kernel_types: List[str] = None,
    output_dim: int = 2048,
    sigma: float = None,
    seed: int = 42,
) -> Dict:
    """
    Test if different kernels produce distinguishable geometry.
    
    If YES: kernel choice = genuine observer difference
    If NO: all kernels see the same geometry (kernel is not an "observer")
    """
    from scipy.linalg import orthogonal_procrustes
    
    if kernel_types is None:
        kernel_types = ['rbf', 'laplacian', 'rq', 'imq', 'matern']
    
    if sigma is None:
        dists = torch.cdist(X[:100], X[:100])
        sigma = float(dists.median())
    
    # Compute embeddings for each kernel
    embeddings = {}
    for kt in kernel_types:
        kernel = create_kernel(kt, X.shape[1], output_dim, sigma, seed)
        embeddings[kt] = kernel.transform(X).cpu().numpy()
    
    # Compute pairwise Procrustes distances
    distances = {}
    ref_kernel = kernel_types[0]
    ref = embeddings[ref_kernel]
    
    print(f"\n{'='*60}")
    print(f"KERNEL DISCRIMINATION: D={output_dim}, σ={sigma:.4f}")
    print(f"{'='*60}")
    print(f"Reference: {ref_kernel}")
    print(f"{'Kernel':<12} {'Procrustes Dist':<15} {'Correlation':<12}")
    print("-" * 40)
    
    for kt in kernel_types:
        emb = embeddings[kt]
        
        # Align to reference
        R, _ = orthogonal_procrustes(emb, ref)
        aligned = emb @ R
        
        # Distance and correlation
        dist = np.linalg.norm(aligned - ref, ord='fro')
        corr = np.corrcoef(aligned.flatten(), ref.flatten())[0, 1]
        
        distances[kt] = {
            'procrustes_distance': float(dist),
            'correlation': float(corr),
        }
        
        print(f"{kt:<12} {dist:<15.4f} {corr:<12.4f}")
    
    # Analysis
    dists_from_ref = [distances[kt]['procrustes_distance'] for kt in kernel_types[1:]]
    mean_dist = np.mean(dists_from_ref)
    
    results = {
        'kernel_types': kernel_types,
        'distances': distances,
        'analysis': {
            'mean_distance_from_ref': float(mean_dist),
            'kernels_distinguishable': mean_dist > 1.0,  # Threshold is data-dependent
            'interpretation': (
                "Kernels produce distinguishable geometry"
                if mean_dist > 1.0 else
                "WARNING: All kernels produce similar geometry"
            ),
        },
    }
    
    print(f"\n{results['analysis']['interpretation']}")
    
    return results


def plot_d_sweep_results(results: Dict, output_path: Path = None):
    """Visualize D-sweep experiment."""
    fig, axes = plt.subplots(1, 2, figsize=(12, 5))
    
    dims = sorted(results['by_dim'].keys())
    corrs = [results['by_dim'][d]['correlation'] for d in dims]
    residuals = [results['by_dim'][d]['mean_residual'] for d in dims]
    normalized = [results['by_dim'][d]['normalized_residual'] for d in dims]
    
    # Plot 1: Approximation quality
    ax1 = axes[0]
    ax1.semilogx(dims, corrs, 'bo-', linewidth=2, markersize=8, base=2)
    ax1.axhline(0.95, color='g', linestyle='--', alpha=0.7, label='Target (0.95)')
    ax1.set_xlabel('Output Dimension D')
    ax1.set_ylabel('Correlation(K, ZZᵀ)')
    ax1.set_title('RFF Approximation Quality')
    ax1.legend()
    ax1.grid(True, alpha=0.3)
    
    # Plot 2: Residuals vs D
    ax2 = axes[1]
    ax2.loglog(dims, residuals, 'ro-', linewidth=2, markersize=8, base=2, label='Raw residual')
    ax2.loglog(dims, normalized, 'g^--', linewidth=2, markersize=8, base=2, label='Normalized (÷√D)')
    
    # Reference line: D^(-0.5) scaling
    ref_line = [residuals[0] * (dims[0] / d) ** 0.5 for d in dims]
    ax2.loglog(dims, ref_line, 'k:', alpha=0.5, label='D⁻⁰·⁵ reference')
    
    ax2.set_xlabel('Output Dimension D')
    ax2.set_ylabel('Procrustes Residual')
    ax2.set_title('Observer Residuals vs D')
    ax2.legend()
    ax2.grid(True, alpha=0.3)
    
    plt.tight_layout()
    
    if output_path:
        plt.savefig(output_path, dpi=150, bbox_inches='tight')
        print(f"\nSaved plot: {output_path}")
    else:
        plt.show()
    
    plt.close()


def main():
    parser = argparse.ArgumentParser(description="RKS Diagnostic Experiments")
    parser.add_argument('--data-path', type=str, help='Path to observer files')
    parser.add_argument('--synthetic', action='store_true', help='Use synthetic data')
    parser.add_argument('--input-dim', type=int, default=24, help='Input dimension for synthetic')
    parser.add_argument('--n-samples', type=int, default=500, help='Number of samples')
    parser.add_argument('--kernel', type=str, default='rbf', help='Kernel type')
    parser.add_argument('--output-dir', type=str, default='diagnostics', help='Output directory')
    parser.add_argument('--all', action='store_true', help='Run all experiments')
    
    args = parser.parse_args()
    
    # Load or generate data
    if args.synthetic or args.data_path is None:
        print(f"Using synthetic data: {args.n_samples} × {args.input_dim}")
        torch.manual_seed(42)
        X = torch.randn(args.n_samples, args.input_dim)
    else:
        X = load_nli_features(Path(args.data_path))
        print(f"Loaded data: {X.shape}")
    
    output_dir = Path(args.output_dir)
    output_dir.mkdir(parents=True, exist_ok=True)
    
    # Run experiments
    results = {}
    
    # 1. D-sweep (always run)
    print("\n" + "=" * 70)
    print("EXPERIMENT 1: D-SWEEP")
    print("=" * 70)
    d_results = run_d_sweep_experiment(X, kernel_type=args.kernel)
    results['d_sweep'] = d_results
    plot_d_sweep_results(d_results, output_dir / 'd_sweep.png')
    
    if args.all:
        # 2. Sigma sweep
        print("\n" + "=" * 70)
        print("EXPERIMENT 2: SIGMA SWEEP")
        print("=" * 70)
        sigma_results = run_sigma_sweep_experiment(X, kernel_type=args.kernel)
        results['sigma_sweep'] = sigma_results
        
        # 3. Kernel discrimination
        print("\n" + "=" * 70)
        print("EXPERIMENT 3: KERNEL DISCRIMINATION")
        print("=" * 70)
        kernel_results = run_kernel_discrimination_experiment(X)
        results['kernel_discrimination'] = kernel_results
    
    # Save results
    results_path = output_dir / 'diagnostic_results.json'
    
    # Convert to JSON-serializable
    def sanitize(obj):
        if isinstance(obj, dict):
            return {k: sanitize(v) for k, v in obj.items()}
        elif isinstance(obj, list):
            return [sanitize(v) for v in obj]
        elif isinstance(obj, (np.integer, np.floating)):
            return float(obj)
        elif isinstance(obj, np.ndarray):
            return obj.tolist()
        return obj
    
    with open(results_path, 'w') as f:
        json.dump(sanitize(results), f, indent=2)
    
    print(f"\n{'='*70}")
    print(f"Results saved to: {results_path}")
    print(f"{'='*70}")
    
    # Final verdict
    print("\n" + "=" * 70)
    print("VERDICT")
    print("=" * 70)
    
    if 'd_sweep' in results:
        analysis = results['d_sweep']['analysis']
        if analysis['residuals_shrinking']:
            print("⚠️  Observer residuals are likely MONTE CARLO NOISE")
            print("    Seeds are not observers - they're just RNG variation")
        else:
            print("✓  Observer residuals PERSIST with increasing D")
            print("    This could indicate genuine observer structure")


if __name__ == "__main__":
    main()

\end{CodeBlock}
\end{document}
